\chapter{曲率时空中的场论与自旋分类}

上一章在平直 Minkowski 时空中建立了从 Poincaré 代数到 Noether 守恒律的完整链条。那里的一切都依赖于一个关键前提：全局惯性坐标系存在，平移生成元 \(P^\mu\) 和 Lorentz 生成元 \(M^{\mu\nu}\) 是全局定义的，普通偏导数 \(\partial_\mu\) 自动协变。一旦时空有了曲率——例如引力场中的弯曲时空或宇宙学背景——这些前提全部失效。没有全局惯性系，平移不再是等距变换，普通偏导数不再协变，Poincaré 群的全局对称性消失。

然而物理并没有崩塌。曲率时空中，全局 Poincaré 对称被两个更局域的结构取代：diffeomorphism covariance（坐标变换的普遍协变性）和局域 Lorentz 对称（每个切空间中独立的 Lorentz 旋转）。上一章构造的 \(\Sigma^{\mu\nu}\)——记录场在 Lorentz 表示中如何旋转的矩阵——在这里获得新的生命：它进入协变导数 \(D_\mu=\partial_\mu-\frac{i}{2}\omega_{\mu ab}\Sigma^{ab}\)，把平直时空的所有场论结构一并推广到曲率时空。不同类型的场（标量、矢量、旋量、张量-旋量）只是 \(\Sigma^{ab}\) 不同，协变导数的统一形式不变。

本章的路线图是这样的。首先解释为什么曲率时空需要 vielbein 和 spin connection，以及协变导数如何统一处理所有 Lorentz 表示。然后构造曲率时空中的作用量，推导 Euler--Lagrange 方程和 Hilbert 能动量张量，并从 diffeomorphism invariance 导出协变守恒律 \(\nabla_\mu T^{\mu\nu}=0\)。接着让度规本身成为动力学变量，得到 Einstein--Hilbert 作用量和 Einstein 方程。最后系统讨论 Wigner 分类：把场的 Lorentz 表示与物理单粒子的 Poincaré 不可约表示联系起来，逐一分析自旋 \(0,\frac{1}{2},1,\frac{3}{2},2\) 的场，以及一般自旋的 Fronsdal 理论，最终到达螺旋度振幅和 little-group scaling——现代量子场论中表示论最直接的应用。

\section{从平直到曲率：vielbein 与局域 Lorentz 框架}

接下来考察为什么不能直接用 \(\partial_\mu\to\nabla_\mu\)。

平直时空中，普通偏导数 \(\partial_\mu\phi\) 对标量场自动协变，对矢量场则需引入 Christoffel 联络得到 \(\nabla_\mu V^\nu=\partial_\mu V^\nu+\Gamma^\nu_{\mu\rho}V^\rho\)。曲率时空中 Christoffel 联络补偿坐标指标的弯曲，这已在前几章的微分几何部分建立。然而对于旋量场，问题出现了：旋量指标不属于任何坐标张量表示，没有 Christoffel 联络可以作用其上。直接写 \(\nabla_\mu\psi=\partial_\mu\psi+\Gamma\psi\) 没有意义——\(\Gamma^\rho_{\mu\nu}\) 的指标是坐标指标，无法作用在旋量指标上。

解决这个困难的物理直觉是：虽然曲率时空中没有全局惯性系，但每一点 \(x\) 的切空间 \(T_xM\) 仍然是 \(d\) 维 Minkowski 矢量空间——切空间中的物理仍然是局域平直的。旋量、矢量等 Lorentz 表示的对象生活在切空间中，而非坐标空间中。因此需要两套指标：坐标指标 \(\mu,\nu\)（被 \(g_{\mu\nu}\) 升降，被 \(\Gamma\) 联络）和切空间指标 \(a,b\)（被 \(\eta_{ab}\) 升降，被 spin connection \(\omega_\mu{}^a{}_b\) 联络）。连接两套指标的桥梁就是 vielbein。

接下来考察Vielbein：切空间与坐标空间的桥梁。

引入 vielbein（tetrad，four-leg）\(e^a{}_\mu(x)\)，它把坐标基矢量 \(e_\mu\) 映射到切空间的 Lorentz 基矢量 \(e_a\)。度规与 Minkowski 度规的关系为
\begin{equation}\label{eq:ms-curved-vielbein}
\boxed{g_{\mu\nu}=e^a{}_\mu\,e^b{}_\nu\,\eta_{ab},}
\end{equation}
逆 vielbein \(e^\mu{}_a\) 满足
\begin{equation}
e^\mu{}_a\,e^a{}_\nu=\delta^\mu{}_\nu,\qquad e^\mu{}_a\,e^b{}_\mu=\delta_a{}^b.
\end{equation}
切空间指标 \(a,b=0,\ldots,d-1\) 用 Minkowski 度规 \(\eta_{ab}\) 升降，坐标指标 \(\mu,\nu\) 用度规 \(g_{\mu\nu}\) 升降。vielbein 把两类指标互相转换：\(V^a=e^a{}_\mu V^\mu\)，\(V^\mu=e^\mu{}_a V^a\)。

vielbein 不是唯一确定的。给定度规 \(g_{\mu\nu}\)，如果 \(e^a{}_\mu\) 是一个合法 vielbein，那么 \(\tilde e^a{}_\mu=\Lambda^a{}_b(x)\,e^b{}_\mu\) 也是——只要 \(\Lambda^a{}_b(x)\) 在每一点 \(x\) 是一个 Lorentz 变换（\(\Lambda^T\eta\Lambda=\eta\)）。这个 \(\Lambda\) 可以依赖于 \(x\)，因为~\eqnrefs{eq:ms-curved-vielbein} 只涉及 \(\eta_{ab}\) 的缩并，局域 Lorentz 旋转在缩并中自动抵消。因此曲率时空具有一个\emph{局域} Lorentz 规范自由度：在每个点独立地旋转切空间基，度规不变。这个局域对称性正是 spin connection 要补偿的对象。

接下来考察Spin connection：为什么需要它。

场 \(\Phi^A\) 若处于某个 Lorentz 表示 \(\Sigma^{ab}\)，在局域 Lorentz 变换下按 \(\Phi\to D(\Lambda(x))\Phi\) 变换。由于 \(\Lambda\) 依赖于 \(x\)，普通导数
\begin{equation*}
\partial_\mu(D\Phi)=D\,\partial_\mu\Phi+(\partial_\mu D)\Phi
\end{equation*}
多出第二项 \((\partial_\mu D)\Phi\)，不再按同一表示协变。物理上，这说"对场求导"和"对场做局域 Lorentz 旋转"这两个操作不交换——导数能"看到"局域 frame 的变化。必须引入一个 connection 来吸收这个额外项，就像 Christoffel 联络吸收坐标基的变化一样。

接下来考察一般表示上的协变导数。

对处于任意 Lorentz 表示 \(\Sigma^{ab}\) 的场，定义协变导数
\begin{equation}\label{eq:ms-curved-Dmu}
\boxed{
D_\mu=\partial_\mu-\frac{i}{2}\omega_{\mu ab}\Sigma^{ab},
}
\end{equation}
其中 \(\omega_{\mu ab}=-\omega_{\mu ba}\) 是 spin connection（取值在 Lorentz 代数中）。要求 \(D_\mu\) 满足协变条件
\begin{equation}\label{eq:ms-curved-D-cov}
D'_\mu\Phi'=D(\Lambda)\,D_\mu\Phi,
\end{equation}
即在变换后的 frame 中对变换后的场求协变导数，等于对原始协变导数做表示变换。这个条件唯一确定了 spin connection 的变换律。

\begin{derivation}[从协变条件导出 spin connection 的变换律]
\emph{思路：}把 \(\Phi'=D\Phi\) 和 \(D'_\mu\) 的定义代入协变条件~\eqnrefs{eq:ms-curved-D-cov}，比较两边，反解出 \(\omega'_\mu\) 与 \(\omega_\mu\) 的关系。计算的关键是把 \(\partial_\mu(D\Phi)\) 用 Leibniz 法则展开，然后与右边的 \(D\,\partial_\mu\Phi\) 抵消。

\emph{第一步：写出协变条件的完整展开。}左边 \(D'_\mu\Phi'=\big(\partial_\mu-\frac{i}{2}\omega'_{\mu ab}\Sigma^{ab}\big)D\Phi\)。把 \(\omega_\mu\equiv-\frac{i}{2}\omega_{\mu ab}\Sigma^{ab}\) 视作 Lie algebra-valued connection，左边简化为
\begin{equation*}
D'_\mu\Phi'=(\partial_\mu D)\Phi+D\,\partial_\mu\Phi+\omega'_\mu\,D\Phi.
\end{equation*}
第一项 \((\partial_\mu D)\Phi\) 来自 Leibniz 法则：\(\partial_\mu\) 作用在乘积 \(D(\Lambda(x))\cdot\Phi(x)\) 上，既对 \(D\) 的 \(x\) 依赖性求导（给出 \((\partial_\mu D)\Phi\)），也对 \(\Phi\) 求导（给出 \(D\,\partial_\mu\Phi\)）。第一项正是需要被 connection 吸收的"额外项"——它测量局域 Lorentz frame 沿 \(x^\mu\) 方向的变化率。

\emph{第二步：展开右边。}右边 \(D\,D_\mu\Phi=D\big(\partial_\mu\Phi+\omega_\mu\Phi\big)\)。利用 \(D\) 是线性算符：
\begin{equation*}
D\,D_\mu\Phi=D\,\partial_\mu\Phi+D\,\omega_\mu\,\Phi.
\end{equation*}
第一项是表示矩阵作用于普通导数，第二项是表示矩阵作用于 connection 乘场。

\emph{第三步：比较两边并消去公共项。}协变条件要求左边 = 右边：
\begin{equation*}
(\partial_\mu D)\Phi+D\,\partial_\mu\Phi+\omega'_\mu\,D\Phi=D\,\partial_\mu\Phi+D\,\omega_\mu\,\Phi.
\end{equation*}
两边都有 \(D\,\partial_\mu\Phi\)，消去后得
\begin{equation*}
(\partial_\mu D)\Phi+\omega'_\mu\,D\Phi=D\,\omega_\mu\,\Phi.
\end{equation*}
这个方程的物理含义：左边的"额外项" \((\partial_\mu D)\Phi\) 加上变换后 connection 的作用 \(\omega'_\mu D\Phi\)，必须等于原始 connection 在表示矩阵下的共轭 \(D\,\omega_\mu\,\Phi\)。

\emph{第四步：解出 \(\omega'_\mu\)。}把 \(\omega'_\mu\) 项移到一边，其余移到另一边：
\begin{equation*}
\omega'_\mu\,D\Phi=D\,\omega_\mu\,\Phi-(\partial_\mu D)\Phi=\big(D\,\omega_\mu-(\partial_\mu D)\big)\Phi.
\end{equation*}
右乘 \(D^{-1}\)（即对表示矩阵取逆），并利用 \(\Phi\) 任意性：
\begin{equation}\label{eq:ms-curved-omega-transform}
\boxed{
\omega'_\mu=D\,\omega_\mu\,D^{-1}-(\partial_\mu D)\,D^{-1}.
}
\end{equation}
写成 Lorentz 矩阵形式即 \(\omega'_\mu=\Lambda\,\omega_\mu\Lambda^{-1}-(\partial_\mu\Lambda)\Lambda^{-1}\)。第一项是共轭作用（表示矩阵的伴随变换），第二项是 inhomogeneous 部分——与电磁规范势的变换律 \(A'_\mu=\Lambda A_\mu\Lambda^{-1}-(\partial_\mu\Lambda)\Lambda^{-1}\) 结构完全相同。这不是巧合：spin connection 正是 Lorentz 规范场的规范势，局域 Lorentz 对称性正是规范群。第二项中的 \((\partial_\mu\Lambda)\Lambda^{-1}\) 补偿了局域 frame 的变化，正如电磁规范势的 inhomogeneous 项补偿规范函数的时空依赖性。

\emph{规范理论类比.}\enspace spin connection $\omega_\mu$ 与 Yang--Mills 规范势 $A_\mu^a T^a$ 的结构完全平行：规范群 $G$（Lorentz $SO(1,3)$ vs 内禀 $SU(N)$），规范势（$\omega_\mu$ vs $A_\mu$），规范变换（$g\in G$ 作用），场强（曲率 $R_{\mu\nu}$ vs $F_{\mu\nu}$），协变导数（$D_\mu=\partial_\mu+\omega_\mu$ vs $D_\mu=\partial_\mu+A_\mu$）。这一类比使引力理论可借用规范理论的所有工具：Wilson 圈（holonomy）、规范固定、BRST 对称、拓扑不变量（Chern 类）。差异在于引力有 tetrad $e_\mu^a$（规范理论无对应物），使引力是"规范理论 + soldering form"，比纯规范理论更丰富。

\emph{Tetrad 形式与 vielbein.}\enspace Tetrad $e_\mu^a$（vielbein，16 个分量）把一般坐标基 $dx^\mu$ 与 Lorentz 标架基 $e^a$ 联系：$e^a=e_\mu^a dx^\mu$。度规 $g_{\mu\nu}=\eta_{ab}e_\mu^a e_\nu^b$ 由 tetrad 构造。Tetrad 的 Lorentz 变换 $e_\mu^a\to\Lambda^a{}_b e_\mu^b$（局域 Lorentz 对称）不改变度规，是规范冗余。spin connection $\omega_\mu^{ab}$ 与 tetrad 通过无挠条件 $T^a=de^a+\omega^a{}_b\wedge e^b=0$（Cartan 第一结构方程）唯一确定：$\omega$ 是 $e$ 的函数（Levi-Civita connection）。这把引力理论表述为 tetrad + spin connection 的 Lorentz 规范理论，是 Einstein--Cartan 形式和圈量子引力的基础。

\emph{Holonomy 与 Wilson 圈.}\enspace spin connection 沿闭合曲线 $C$ 的路径有序指数 $W(C)=\mathcal{P}\exp\oint_C\omega_\mu dx^\mu$ 是 Lorentz holonomy，描述标架沿 $C$ 平移回来的旋转。对无穷小回路，$W\approx1+\frac12 R_{\mu\nu}\sigma^{\mu\nu}$（$\sigma^{\mu\nu}$ 为回路面积），曲率是 holonomy 的无穷小生成元。Holonomy 的共轭类（Lorentz 转角）是微分同胚不变量，与 Wilson 圈在规范理论中的角色完全平行。在圈量子引力中，holonomy 是基本变量，与标架的体积元共同给出量子化的几何算符（面积、体积本征值为离散谱），是引力量子化的背景无关方案。
\end{derivation}

接下来考察不同表示只是 \(\Sigma^{ab}\) 不同。

协变导数~\eqnrefs{eq:ms-curved-Dmu} 的统一性是整个框架的核心力量。同一个公式适用于所有 Lorentz 表示，只是 \(\Sigma^{ab}\) 随场类型变化：

\begin{itemize}
\item \textbf{标量}：\(\Sigma^{ab}=0\)，无内禀旋转，协变导数退化为普通导数 \(D_\mu\phi=\partial_\mu\phi\)。
\item \textbf{Lorentz 矢量}：\((\Sigma^{ab})^c{}_d=i(\eta^{ac}\delta^b{}_d-\eta^{bc}\delta^a{}_d)\)，代入后 \(D_\mu V^a=\partial_\mu V^a+\omega_\mu{}^a{}_b V^b\)——与 Christoffel 联络作用在坐标矢量上形式一致。
\item \textbf{Dirac 旋量}：\(\Sigma^{ab}=\frac{i}{4}[\gamma^a,\gamma^b]\)（由上一章~\eqnrefs{eq:ms-poincare-Sigma-dirac} 给出），协变导数 \(D_\mu\psi=\big(\partial_\mu-\frac{i}{2}\omega_{\mu ab}\Sigma^{ab}\big)\psi\)。
\item \textbf{任意 tensor-spinor}：总表示生成元是各指标贡献之和，\(\Sigma^{ab}_{\text{total}}=\Sigma^{ab}_{\text{tensor}}+\Sigma^{ab}_{\text{spinor}}\)。
\end{itemize}

这种统一性意味着：无论场多复杂，只需识别它的 Lorentz 表示，代入正确的 \(\Sigma^{ab}\)，协变导数自动给出正确的联络项。不需要为每种场类型单独推导。

接下来考察Tetrad postulate：联络的一致性。

矢量场可以同时用切空间指标 \(V^a\) 和坐标指标 \(V^\mu=e^\mu{}_a V^a\) 描述。两种描述必须给出相同的协变导数。要求 vielbein 本身协变导数为零（\(D_\mu e^a{}_\nu=0\)），即"vielbein 是平移的"，给出
\begin{equation}\label{eq:ms-curved-tetrad}
\boxed{
\partial_\mu e^a{}_\nu+\omega_\mu{}^a{}_b\,e^b{}_\nu-\Gamma^\rho_{\mu\nu}\,e^a{}_\rho=0.
}
\end{equation}
这个 tetrad postulate 把 spin connection \(\omega_\mu{}^a{}_b\) 与 affine connection \(\Gamma^\rho_{\mu\nu}\) 联系起来：给定 vielbein 和 \(\Gamma\)，spin connection 被唯一确定。若 \(\Gamma\) 取 Levi--Civita 联络（无挠、度规相容），则 \(\omega_\mu{}^a{}_b\) 是对应的 Levi--Civita spin connection。tetrad postulate 保证了无论用切空间还是坐标指标计算，物理结果一致。

\section{曲率时空中的作用量与 Euler--Lagrange 方程}

有了协变导数，构造作用量只需把平直时空公式中的 \(\eta_{\mu\nu}\to g_{\mu\nu}\)、\(\partial_\mu\to D_\mu\)、\(d^dx\to d^dx\sqrt{|g|}\)。一般形式为
\begin{equation}\label{eq:ms-curved-action}
\boxed{
S[\Phi;g]=\int d^dx\,\sqrt{|g|}\;\mathcal{L}(\Phi^A,D_\mu\Phi^A,g_{\mu\nu},R,R_{\mu\nu},R_{\mu\nu\rho\sigma},\ldots),
}
\end{equation}
其中拉氏密度可以依赖场、协变导数、度规以及曲率张量（非最小耦合）。最简单的最小耦合情形只含 \(\Phi\) 和 \(D_\mu\Phi\)。

\begin{derivation}[曲率时空 Euler--Lagrange 方程的推导]
\emph{思路：}与平直情形~\eqnrefs{eq:ms-poincare-EL} 完全平行，但用协变分部积分代替普通分部积分。关键区别是：\(\sqrt{|g|}\) 因子使测度不再是平坦的，\(\nabla_\mu\) 代替 \(\partial_\mu\) 使分部积分多出联络项——但两者恰好抵消，最终形式与平直情形结构相同。

\emph{第一步：写出变分。}场 \(\Phi^A\to\Phi^A+\delta\Phi^A\) 引起
\begin{equation*}
\delta S=\int d^dx\sqrt{|g|}\left[\frac{\partial\mathcal{L}}{\partial\Phi^A}\delta\Phi^A+\frac{\partial\mathcal{L}}{\partial(D_\mu\Phi^A)}D_\mu\delta\Phi^A\right].
\end{equation*}
第一项来自 \(\mathcal{L}\) 对场的显式依赖，第二项来自 \(\mathcal{L}\) 对协变导数的依赖。

\emph{第二步：定义共轭动量。}为简洁起见定义
\begin{equation*}
\Pi_A^\mu\equiv\frac{\partial\mathcal{L}}{\partial(D_\mu\Phi^A)}.
\end{equation*}

\emph{第三步：协变分部积分。}关键恒等式：对任意矢量场 \(V^\mu\) 和标量密度 \(\sqrt{|g|}\)，有
\begin{equation*}
\int d^dx\sqrt{|g|}\,\nabla_\mu V^\mu=\int d^dx\,\partial_\mu(\sqrt{|g|}V^\mu).
\end{equation*}
右端是全散度，积分后化为边界项。因此协变分部积分（忽略边界）为
\begin{equation*}
\int d^dx\sqrt{|g|}\,\Pi_A^\mu\,D_\mu\delta\Phi^A=-\int d^dx\sqrt{|g|}\,(D_\mu\Pi_A^\mu)\delta\Phi^A.
\end{equation*}
这里 \(D_\mu\delta\Phi^A\) 中 \(\delta\Phi^A\) 与 \(\Phi^A\) 处于同一表示，因此 \(D_\mu\) 的联络项在分部积分中正确传递。

\emph{第四步：代入并读出方程。}
\begin{equation*}
\delta S=\int d^dx\sqrt{|g|}\left[\frac{\partial\mathcal{L}}{\partial\Phi^A}-D_\mu\Pi_A^\mu\right]\delta\Phi^A+\text{(boundary)}.
\end{equation*}
边界项不贡献（变分在边界消失）。由 \(\delta\Phi^A\) 任意性，得
\begin{equation}\label{eq:ms-curved-EL}
\boxed{
\frac{\partial\mathcal{L}}{\partial\Phi^A}-D_\mu\frac{\partial\mathcal{L}}{\partial(D_\mu\Phi^A)}=0.
}
\end{equation}
形式与平直 Euler--Lagrange 方程~\eqnrefs{eq:ms-poincare-EL} 完全一致，只是 \(\partial_\mu\to D_\mu\)。这正是协变导数统一构造的好处：平直时空的所有推导可以"形式地"推广，只需把偏导数换成协变导数。

\emph{最小耦合原理.}\enspace 平直时空的 Lagrangian $\mathcal{L}_{\text{flat}}(\Phi,\partial\Phi)$ 推广到弯曲时空的"最小耦合"规则：(1) $\eta_{\mu\nu}\to g_{\mu\nu}$（Minkowski 度规换一般度规）；(2) $\partial_\mu\to\nabla_\mu$（偏导数换协变导数）；(3) $d^4x\to\sqrt{-g}\,d^4x$（平坦测度换不变测度）。例如标量场 $\mathcal{L}=\frac12\eta^{\mu\nu}\partial_\mu\phi\partial_\nu\phi-\frac12m^2\phi^2$ 推广为 $\mathcal{L}=\frac12g^{\mu\nu}\nabla_\mu\phi\nabla_\nu\phi-\frac12m^2\phi^2=\frac12g^{\mu\nu}\partial_\mu\phi\partial_\nu\phi-\frac12m^2\phi^2$（标量场 $\nabla_\mu\phi=\partial_\mu\phi$）。最小耦合不是唯一推广（可加曲率耦合 $\xi R\phi^2$），但最简单且在低能极限下普适。

\emph{标量场的弯曲时空方程.}\enspace 标量场的作用量 $S=\int\left(\frac12g^{\mu\nu}\partial_\mu\phi\partial_\nu\phi-\frac12m^2\phi^2\right)\sqrt{-g}\,d^4x$ 的 Euler--Lagrange 方程为 $\frac{1}{\sqrt{-g}}\partial_\mu(\sqrt{-g}\,g^{\mu\nu}\partial_\nu\phi)+m^2\phi=0$，即 $(\Box_g+m^2)\phi=0$（$\Box_g=\frac{1}{\sqrt{-g}}\partial_\mu(\sqrt{-g}\,g^{\mu\nu}\partial_\nu)$ 为弯曲时空 d'Alembert 算符）。在 de Sitter 时空中（$R=12H^2$），共形耦合标量场 $(\Box_g+m^2+\xi R)\phi=0$（$\xi=1/6$ 为共形耦合）的解与平直时空解通过共形变换关联，是宇宙学暴胀中曲率微扰的理论基础。

\emph{Maxwell 场的弯曲时空推广.}\enspace Maxwell 作用量 $S=-\frac14\int F_{\mu\nu}F^{\mu\nu}\sqrt{-g}\,d^4x$（$F_{\mu\nu}=\nabla_\mu A_\nu-\nabla_\nu A_\mu=\partial_\mu A_\nu-\partial_\nu A_\mu$，协变导数项抵消）给出弯曲时空 Maxwell 方程 $\nabla_\mu F^{\mu\nu}=0$（$\nabla_\mu$ 含 Christoffel 符号）。在弯曲时空中，光线的弯曲（引力透镜）和红移（宇宙学红移）直接来自 Maxwell 方程在弯曲度规上的解。规范不变性 $A_\mu\to A_\mu+\nabla_\mu\alpha$（等价于 $A_\mu\to A_\mu+\partial_\mu\alpha$，因 $\nabla_\mu\nabla_\nu\alpha-\nabla_\nu\nabla_\mu\alpha=0$ 对标量 $\alpha$）在弯曲时空中保持，是电磁规范对称的协变推广。
\end{derivation}

\section{Hilbert 能动量张量与协变守恒}

接下来考察Hilbert 定义。

平直时空中，能动量张量从平移对称性的 Noether 流导出（canonical 张量~\eqnrefs{eq:ms-poincare-T-can}）。曲率时空中没有全局平移对称性，但有一个更直接的替代方案：把度规 \(g_{\mu\nu}\) 视为背景场，定义 Hilbert 能动量张量为作用量对度规的变分
\begin{equation}\label{eq:ms-curved-T-hilbert}
\boxed{
T_{\mu\nu}=-\frac{2}{\sqrt{|g|}}\frac{\delta S_m}{\delta g^{\mu\nu}}.
}
\end{equation}
因为 \(\delta g^{\mu\nu}=\delta g^{\nu\mu}\)，所以 \(T_{\mu\nu}=T_{\nu\mu}\) 自动对称——不需要 Belinfante 改进。这是 Hilbert 定义相对于 canonical 定义的一大优势：对称性是定义自动保证的，而非需要额外修正的。

接下来考察Diffeomorphism invariance 推出协变守恒。

曲率时空中没有全局 Poincaré 对称，但作用量~\eqnrefs{eq:ms-curved-action} 满足 diffeomorphism covariance：在任何坐标变换下形式不变。无穷小 diffeomorphism \(x^\mu\to x^\mu+\xi^\mu(x)\) 生成度规变分
\begin{equation*}
\delta_\xi g_{\mu\nu}=\nabla_\mu\xi_\nu+\nabla_\nu\xi_\mu,
\end{equation*}
即度规的 Lie 导数。在物质场方程成立时（on-shell），作用量变分为零——这就是 diffeomorphism invariance 的 Noether 恒等式。

\begin{derivation}[从 diffeomorphism invariance 导出协变守恒律]
\emph{思路：}把无穷小 diffeomorphism 引起的度规变分代入 Hilbert 张量定义，利用 \(T_{\mu\nu}\) 对称性简化，然后协变分部积分提出任意矢量 \(\xi^\nu\)，其系数必须为零。核心逻辑链：diffeomorphism invariance \(\to\) on-shell \(\delta_\xi S=0\) \(\to\) \(T^{\mu\nu}\) 恒等式 \(\to\) 协变守恒。

\emph{第一步：度规变分与作用量变分。}无穷小 diffeomorphism \(x^\mu\to x^\mu+\xi^\mu(x)\) 使度规产生 Lie 导数 \(\delta_\xi g_{\mu\nu}=\nabla_\mu\xi_\nu+\nabla_\nu\xi_\mu\)。由 Hilbert 定义~\eqnrefs{eq:ms-curved-T-hilbert}，\(\delta_g S_m=-\frac{1}{2}\int d^dx\sqrt{|g|}\,T^{\mu\nu}\delta g_{\mu\nu}\)。代入：
\begin{equation*}
\delta_\xi S_m=-\frac{1}{2}\int d^dx\sqrt{|g|}\,T^{\mu\nu}(\nabla_\mu\xi_\nu+\nabla_\nu\xi_\mu).
\end{equation*}

\emph{第二步：利用 \(T^{\mu\nu}\) 对称性。}Hilbert 张量自动对称（\(T^{\mu\nu}=T^{\nu\mu}\)，因为 \(\delta g^{\mu\nu}=\delta g^{\nu\mu}\)）。因此两项贡献相等：
\begin{equation*}
T^{\mu\nu}\nabla_\mu\xi_\nu+T^{\mu\nu}\nabla_\nu\xi_\mu=2\,T^{\mu\nu}\nabla_\mu\xi_\nu.
\end{equation*}
验证：第二项交换哑指标 \(\mu\leftrightarrow\nu\) 得 \(T^{\nu\mu}\nabla_\mu\xi_\nu\)，利用 \(T^{\nu\mu}=T^{\mu\nu}\) 后与第一项相同。因此
\begin{equation*}
\delta_\xi S_m=-\int d^dx\sqrt{|g|}\,T^{\mu\nu}\nabla_\mu\xi_\nu.
\end{equation*}

\emph{第三步：协变分部积分。}利用 Leibniz 法则把协变导数从 \(\xi_\nu\) 移到 \(T^{\mu\nu}\) 上：
\begin{equation*}
T^{\mu\nu}\nabla_\mu\xi_\nu=\nabla_\mu(T^{\mu\nu}\xi_\nu)-(\nabla_\mu T^{\mu\nu})\xi_\nu.
\end{equation*}
第一项是全协变散度 \(\nabla_\mu V^\mu\)（其中 \(V^\mu=T^{\mu\nu}\xi_\nu\)）。关键恒等式：\(\int d^dx\sqrt{|g|}\,\nabla_\mu V^\mu=\int d^dx\,\partial_\mu(\sqrt{|g|}V^\mu)\)，右端是普通全散度，积分后化为边界项。由于 \(\xi^\nu\) 在边界上消失（或场在无穷远衰减足够快），边界项不贡献。剩下
\begin{equation*}
\delta_\xi S_m=\int d^dx\sqrt{|g|}\,(\nabla_\mu T^{\mu\nu})\xi_\nu.
\end{equation*}

\emph{第四步：利用 on-shell 条件和 \(\xi^\nu\) 任意性。}在物质场方程成立时（on-shell），diffeomorphism invariance 给出 \(\delta_\xi S_m=0\)。由于 \(\xi^\nu(x)\) 是任意矢量场——它在每一点可以取任意方向和大小——要使积分对所有 \(\xi^\nu\) 为零，被积函数的系数必须逐点为零：
\begin{equation}\label{eq:ms-curved-T-conservation}
\boxed{\nabla_\mu T^{\mu\nu}=0.}
\end{equation}
这就是协变守恒律。与平直时空的 \(\partial_\mu T^{\mu\nu}=0\) 相比，普通导数被协变导数取代——曲率可以通过联络项产生或吸收能量-动量流。

\emph{微分同胚不变性与守恒律.}\enspace 弯曲时空的 $\nabla_\mu T^{\mu\nu}=0$ 来自作用量的微分同胚不变性（一般坐标变换下的不变性），而非平移对称性。微分同胚 $\xi^\mu(x)$ 生成坐标变换 $x^\mu\to x^\mu+\xi^\mu$，作用量不变给出 $\nabla_\mu T^{\mu\nu}=0$。这与 Noether 定理一致：对称性（微分同胚）给出守恒律（协变守恒）。但与平直时空不同，$\nabla_\mu T^{\mu\nu}=0$ 不能直接积分为守恒荷（因协变导数含联络项），需要额外结构（Killing 矢量）才能构造守恒量。

\emph{Killing 矢量与守恒荷.}\enspace 若度规有 Killing 矢量 $K^\mu$（$\nabla_\mu K_\nu+\nabla_\nu K_\mu=0$），则 $J_K=\int_\Sigma T^{\mu\nu}K_\nu\,d\Sigma_\mu$ 是守恒荷：$\partial_\mu(T^{\mu\nu}K_\nu)=(\nabla_\mu T^{\mu\nu})K_\nu+T^{\mu\nu}\nabla_\mu K_\nu=0$（第一项由协变守恒，第二项由 Killing 方程和 $T^{\mu\nu}$ 对称性）。Minkowski 时空有 10 个 Killing 矢量（4 平移 + 6 Lorentz），给出 10 个守恒荷（能量-动量 + 角动量）。Schwarzschild 时空有 2 个（时间平移 $\partial_t$ 和轴向旋转 $\partial_\phi$），给出能量和角动量守恒。一般时空无 Killing 矢量，无全局守恒荷。

\emph{Landau--Lifshitz 伪张量.}\enspace $\nabla_\mu T^{\mu\nu}=0$ 不是真正的局部守恒律（不能积分为不依赖于体积边界的荷）。构造真正的守恒律需引入 Landau--Lifshitz 伪张量 $t^{\mu\nu}_{\text{LL}}$：$\partial_\mu(-g)(T^{\mu\nu}+t^{\mu\nu}_{\text{LL}})=0$，其中 $t^{\mu\nu}_{\text{LL}}$ 含度规导数的非线性项（引力场自身的能量-动量）。伪张量不是张量（坐标变换下不协变），但给出引力波能量传输和 ADM 贪的定义。$t^{\mu\nu}_{\text{LL}}$ 的非张量性反映引力能量不能局部化（等效原理：局部可消去引力场），只能在渐近平直区域定义总能量（ADM 贪）。

\emph{引力波能量传输.}\enspace 在渐近平直时空中，引力波在远处携带的能量由 Landau--Lifshitz 伪张量的 $r\to\infty$ 极限给出：$\frac{dE}{dt}=-\lim_{r\to\infty}\int t^{0r}_{\text{LL}}\,r^2d\Omega$。对平面引力波 $h_{ij}^{TT}(t-z)$，能量通量 $t^{0z}_{\text{LL}}=\frac{1}{32\pi G}\langle\dot h_{ij}^{TT}\dot h_{ij}^{TT}\rangle$（时间平均），与电磁波 $T^{0z}=\frac{1}{4\pi}|\dot{\vb A}|^2$ 结构平行。引力波能量传输的实验验证（LIGO 2015，双黑洞合并 GW150914）直接观测到 $t^{\mu\nu}_{\text{LL}}$ 的物理效应，是广义相对论非线性效应的首次直接验证。

\emph{ADM 贪与正质量定理.}\enspace 在渐近平直时空（$g_{\mu\nu}\to\eta_{\mu\nu}+O(1/r)$），ADM 贪 $M_{\text{ADM}}=\frac{1}{16\pi G}\lim_{r\to\infty}\oint(n^i\partial_j g_{ij}-n^j\partial_i g_{ij})\,r^2d\Omega$ 是 Landau--Lifshitz 伪张量的全局积分。正质量定理（Schoen--Yau 1979，Witten 1981）证明 $M_{\text{ADM}}\ge0$（等号当且仅当 Minkowski），是广义相对论的稳定性结果。ADM 贪在圈量子引力中给出面积量子化（Ashtekar 变量），在弦论中匹配弦质量谱，是引力能量在全局层面的精确定义。
\end{derivation}

这就是曲率时空中的协变守恒律。与平直时空的 \(\partial_\mu T^{\mu\nu}=0\)（~\eqnrefs{eq:ms-poincare-T-conservation}）相比，普通导数被协变导数取代。物理含义：能量-动量不能在弯曲时空中"自由流动"——曲率本身可以产生或吸收能量-动量的流。只有在平直极限（\(\Gamma=0\)）下，\(\nabla_\mu\to\partial_\mu\)，才恢复普通守恒。

接下来考察Killing 矢量与守恒荷。

协变守恒律 \(\nabla_\mu T^{\mu\nu}=0\) 本身不直接给出守恒荷——因为 \(\nabla_\mu\) 含联络项，\(T^{0\nu}\) 的空间积分一般不守恒。但如果时空存在 Killing 矢量 \(\xi^\mu\)（满足 Killing 方程 \(\nabla_{(\mu}\xi_{\nu)}=0\)），则可以构造守恒流
\begin{equation}\label{eq:ms-curved-killing-current}
\boxed{
J^\mu_\xi=T^\mu{}_\nu\xi^\nu,\qquad \nabla_\mu J^\mu_\xi=0.
}
\end{equation}
Killing 矢量代表时空的等距变换——沿 \(\xi^\mu\) 方向的"运动"保持度规不变。物理上，每个 Killing 矢量对应一个残余对称性：虽然全局 Poincaré 对称消失了，但时空仍可能保留部分等距（如 Schwarzschild 时空的静态性和球对称性）。对应的守恒荷
\begin{equation}\label{eq:ms-curved-killing-charge}
Q_\xi=\int_\Sigma d\Sigma_\mu\,T^\mu{}_\nu\xi^\nu
\end{equation}
不随时间变化。例如：静态时空的时间平移 Killing 矢量给出能量守恒，球对称时空的三个旋转 Killing 矢量给出角动量守恒。

\section{曲率时空中的典型场}

现在把一般框架应用到具体场类型，验证协变导数的统一性。

接下来考察标量场。

标量场最简单：\(\Sigma^{ab}=0\)，协变导数退化为普通导数。含非最小曲率耦合的作用量为
\begin{equation}
S_\phi=-\frac{1}{2}\int d^dx\sqrt{|g|}\,\big[\nabla_\mu\phi\nabla^\mu\phi+m^2\phi^2+\xi R\phi^2\big],
\end{equation}
其中 \(\xi\) 是非最小耦合参数。\(\xi=0\) 是最小耦合，\(\xi=\frac{d-2}{4(d-1)}\)（四维时 \(\xi=\frac{1}{6}\)）是共形耦合——后者使标量场在共形变换下具有良好的协变性。变分给出
\begin{equation}\label{eq:ms-curved-scalar-eom}
\boxed{(\Box_g-m^2-\xi R)\phi=0,}
\end{equation}
其中 \(\Box_g=g^{\mu\nu}\nabla_\mu\nabla_\nu\) 是曲率时空的 d'Alembert 算符。与平直 Klein--Gordon 方程 \((\Box-m^2)\phi=0\) 相比，多了曲率耦合项 \(\xi R\phi\)——当 \(\xi\neq0\) 时，曲率直接贡献有效质量 \(m_{\text{eff}}^2=m^2+\xi R\)。

接下来考察Maxwell 场。

电磁势 \(A_\mu\) 是坐标矢量，协变导数作用其上。但场强 \(F_{\mu\nu}=\nabla_\mu A_\nu-\nabla_\nu A_\mu\) 中 Christoffel 项恰好消去：
\begin{equation*}
F_{\mu\nu}=\partial_\mu A_\nu-\Gamma^\rho_{\mu\nu}A_\rho-\partial_\nu A_\mu+\Gamma^\rho_{\nu\mu}A_\rho=\partial_\mu A_\nu-\partial_\nu A_\mu,
\end{equation*}
因为 Levi--Civita 联络无挠（\(\Gamma^\rho_{\mu\nu}=\Gamma^\rho_{\nu\mu}\)）。因此场强在曲率时空中形式不变。作用量
\begin{equation}
S_A=-\frac{1}{4}\int d^dx\sqrt{|g|}\,F_{\mu\nu}F^{\mu\nu}
\end{equation}
给出运动方程 \(\nabla_\mu F^{\mu\nu}=0\)——这里 \(\nabla_\mu\) 不可省略，因为 \(F^{\mu\nu}\) 带指标。Hilbert 能动量张量为
\begin{equation}
T_{\mu\nu}=F_{\mu\rho}F_\nu{}^\rho-\frac{1}{4}g_{\mu\nu}F_{\rho\sigma}F^{\rho\sigma},
\end{equation}
与平直时空形式相同，但指标升降用 \(g_{\mu\nu}\) 而非 \(\eta_{\mu\nu}\)。

接下来考察Dirac 场。

旋量场是协变导数统一性最显著的体现。弯曲时空 gamma 矩阵定义为
\begin{equation}
\gamma^\mu(x)=e^\mu{}_a(x)\gamma^a,
\end{equation}
它们满足弯曲 Clifford 代数
\begin{equation}
\{\gamma^\mu,\gamma^\nu\}=2g^{\mu\nu}.
\end{equation}
旋量协变导数 \(D_\mu\psi=\big(\partial_\mu-\frac{i}{2}\omega_{\mu ab}\Sigma^{ab}\big)\psi\) 中，\(\Sigma^{ab}=\frac{i}{4}[\gamma^a,\gamma^b]\) 是上一章~\eqnrefs{eq:ms-poincare-Sigma-dirac} 给出的旋量表示生成元。Hermitian Dirac 作用量为
\begin{equation}\label{eq:ms-curved-dirac-action}
\boxed{
S_D=\int d^dx\,e\left[\frac{i}{2}\big(\bar\psi\gamma^\mu D_\mu\psi-(D_\mu\bar\psi)\gamma^\mu\psi\big)-m\bar\psi\psi\right],
}
\end{equation}
其中 \(e\equiv\det(e^a{}_\mu)=\sqrt{|g|}\)。对称形式 \(\frac{i}{2}(\bar\psi\gamma^\mu D_\mu\psi-(D_\mu\bar\psi)\gamma^\mu\psi)\) 保证作用量实值。对 \(\bar\psi\) 变分给出弯曲时空 Dirac 方程
\begin{equation}
(i\gamma^\mu D_\mu-m)\psi=0.
\end{equation}
这个方程在形式上与平直 Dirac 方程 \((i\gamma^\mu\partial_\mu-m)\psi=0\) 完全一致，只是 \(\partial_\mu\to D_\mu\)、\(\gamma^\mu\) 变为位置依赖。这正是协变导数构造的力量：平直时空的场方程可以"机械地"推广，协变导数自动处理所有曲率效应。

\section{动态几何：Einstein--Hilbert 作用量}

到目前为止，度规 \(g_{\mu\nu}\) 一直是外部背景场——给定时空几何，研究物质场在其上的动力学。现在迈出关键一步：让度规本身成为动力学变量。物理上，这意味着引力场被量子化（至少在经典层面，度规的运动方程从作用量导出）。

最简单的度规作用量是 Einstein--Hilbert 作用量
\begin{equation}\label{eq:ms-curved-EH}
S_{\text{EH}}[g]=\frac{1}{2\kappa}\int d^dx\sqrt{|g|}\,(R-2\Lambda),
\end{equation}
其中 \(R\) 是标量曲率，\(\Lambda\) 是宇宙学常数，\(\kappa=8\pi G\) 是引力耦合常数（\(G\) 是 Newton 常数）。总作用量为
\begin{equation}
S[g,\Phi]=S_{\text{EH}}[g]+S_m[g,\Phi].
\end{equation}
对 \(g^{\mu\nu}\) 变分：Einstein--Hilbert 部分给出 Einstein 张量 \(G_{\mu\nu}=R_{\mu\nu}-\frac{1}{2}g_{\mu\nu}R\)（加上宇宙学项 \(\Lambda g_{\mu\nu}\)），物质部分给出 Hilbert 能动量张量~\eqnrefs{eq:ms-curved-T-hilbert}。合并得 Einstein 方程
\begin{equation}\label{eq:ms-curved-einstein}
\boxed{G_{\mu\nu}+\Lambda g_{\mu\nu}=\kappa T_{\mu\nu}.}
\end{equation}

Einstein 方程有一个深刻的自洽性。Bianchi 恒等式
\begin{equation}
\nabla_\mu G^{\mu\nu}\equiv 0
\end{equation}
是几何恒等式——对任何度规都成立，不依赖场方程。对~\eqnrefs{eq:ms-curved-einstein} 取协变散度：
\begin{equation*}
\nabla_\mu(G^{\mu\nu}+\Lambda g^{\mu\nu})=\kappa\,\nabla_\mu T^{\mu\nu}.
\end{equation*}
左边：\(\nabla_\mu G^{\mu\nu}=0\)（Bianchi），\(\nabla_\mu(\Lambda g^{\mu\nu})=0\)（度规相容 \(\nabla_\mu g^{\mu\nu}=0\)）。因此左边为零，推出
\begin{equation}
\nabla_\mu T^{\mu\nu}=0.
\end{equation}
物质能动量的协变守恒不是独立假设，而是 Einstein 方程的自动推论。几何方程的恒等式（Bianchi）保证了物质守恒——这是广义相对论最优美的自洽性之一。物理上，引力场方程决定时空如何弯曲，弯曲时空的几何恒等式反过来约束物质如何运动。

\section{Wigner 分类：从场表示到物理粒子}

接下来考察场表示与粒子表示的区别。

到此必须强调一个容易混淆的根本区别：场的 Lorentz 表示\(\neq\)单粒子的 Poincaré 不可约表示。

场 \(\Phi^A(x)\) 的指标 \(A\) 属于有限维 Lorentz 表示，描述场在 Lorentz 变换下的变换律。例如矢量场 \(A_\mu\) 有 4 个 Lorentz 分量，属于 \((\frac{1}{2},\frac{1}{2})\) 表示。但物理单粒子态是 Hilbert 空间中的向量，它们按照 Poincaré 群的\emph{无穷维}酉表示分类——这是 Wigner 1939 年的分类方案。无质量光子只有 \(h=\pm1\) 两个螺旋度态，而非 4 个——场方程和规范对称性把 4 个 Lorentz 分量投影到 2 个物理态。

Wigner 分类的关键是两个 Casimir 不变量——与所有 Poincaré 生成元对易，因此在不可约表示上是常数。第一 Casimir 是质量平方
\begin{equation}
C_1=P^2=P_\mu P^\mu,
\end{equation}
对质量为 \(m\) 的粒子给出 \(P^2=-m^2\)（"mostly plus"度规下）。第二 Casimir 由 Pauli--Lubanski 向量 \(W^\mu=\frac{1}{2}\epsilon^{\mu\nu\rho\sigma}P_\nu M_{\rho\sigma}\) 构造：
\begin{equation}
C_2=W^2=W_\mu W^\mu.
\end{equation}
两个 Casimir 的值完全标记 Poincaré 群的不可约酉表示——即物理粒子类型。

接下来考察Massive 粒子：little group \(SO(3)\)。

四维 massive 粒子满足 \(P^2=-m^2\)，\(m>0\)。可以进入静止系 \(p^\mu=(m,\mathbf{0})\)。保持该动量不变的 Lorentz 变换组成 little group（Wigner 术语）\(SO(3)\)（双覆盖 \(SU(2)\)）——就是普通三维旋转群。因此 massive 单粒子态由 \(SU(2)\) 的不可约表示分类，即由自旋量子数
\begin{equation}
s=0,\frac{1}{2},1,\frac{3}{2},\ldots
\end{equation}
标记，具有 \(2s+1\) 个自旋极化态。Pauli--Lubanski Casimir 在该表示上的值为
\begin{equation}\label{eq:ms-curved-W2-massive}
\boxed{W^2=m^2 s(s+1).}
\end{equation}
因此 massive spin 是真正的 Lorentz invariant 表示标签——不同自旋的粒子属于不同的 Poincaré 不可约表示。

接下来考察Massless 粒子：little group \(ISO(2)\)。

无质量粒子满足 \(P^2=0\)，不能进入静止系。取标准动量 \(k^\mu=(E,0,0,E)\)（沿 \(z\) 方向传播）。保持它的 little group 是 \(ISO(2)\)——二维 Euclid 群，由绕 \(z\) 轴的旋转 \(SO(2)\) 和两个横向"平移"生成。在普通局域量子场论中，取平移子群平凡作用（"finite helicity" 表示），此时 Pauli--Lubanski 向量满足
\begin{equation}\label{eq:ms-curved-W-helicity}
\boxed{W^\mu=h\,P^\mu,}
\end{equation}
其中 \(h\) 就是 helicity（螺旋度）。对 proper orthochronous Lorentz group，massless helicity 是 Lorentz invariant。一个 parity-invariant 的实无质量玻色场通常同时含有 \(h=+s\) 和 \(h=-s\) 两个螺旋度（手征理论可以只保留一侧）。

接下来考察自由度对比。

massive 与 massless 的自由度差异是量子场论中最基本的区分之一：

\begin{center}
\begin{tabular}{lll}
\hline
 & massive spin-\(s\) & massless helicity-\(s\) \\
\hline
自由度 & \(2s+1\) & \(2\)（\(h=\pm s\)） \\
little group & \(SO(3)\) & \(ISO(2)\) \\
Casimir & \(W^2=m^2s(s+1)\) & \(W^\mu=hP^\mu\) \\
\hline
\end{tabular}
\end{center}

同一"spin"标签下，massive 粒子有 \(2s+1\) 个极化态，massless 粒子只有 2 个（手征场甚至只有 1 个）。这个差异的根源是 little group 的不同：\(SO(3)\) 的不可约表示由自旋 \(s\) 标记（维数 \(2s+1\)），而 \(ISO(2)\) 的 finite helicity 表示由单一 helicity \(h\) 标记（维数 1，parity 配对后维数 2）。

\section{自旋 0：标量场}

最简单的自旋分类。标量场属于 Lorentz 表示 \((0,0)\)，\(\Sigma^{\mu\nu}=0\)，无内禀旋转。作用量
\begin{equation}
\mathcal{L}=-\frac{1}{2}(\partial\phi)^2-\frac{1}{2}m^2\phi^2
\end{equation}
给出 Klein--Gordon 方程 \((\Box-m^2)\phi=0\)，物理自由度为 1——无自旋极化。

量子场论中标量场的常见应用：Higgs 标量（势 \(V(\phi)=\frac{1}{2}m^2\phi^2+\frac{\lambda}{4!}\phi^4\)）、Yukawa 耦合 \(-y\phi\bar\psi\psi\)（标量与费米子的最基本耦合）、Goldstone boson（导数耦合的有效场论 \(\mathcal{L}_{\text{EFT}}=\frac{1}{2}(\partial\pi)^2+\frac{c_4}{\Lambda^4}(\partial\pi)^4+\cdots\)）、axion（赝标量，耦合 \(-iy_5\phi\bar\psi\gamma^5\psi\)）、inflaton（宇宙学暴胀标量场）。

\section{自旋 1/2：Weyl、Dirac 与 Majorana}

四维中 Lorentz 群的双覆盖 \(Spin(1,3)\simeq SL(2,\mathbb{C})\)，有限维不可约表示由 \((j_L,j_R)\) 标记。左手 Weyl 旋量 \(\chi_\alpha\in(\frac{1}{2},0)\)，右手 Weyl 旋量 \(\bar\eta^{\dot\alpha}\in(0,\frac{1}{2})\)，Dirac 旋量 \(\psi\in(\frac{1}{2},0)\oplus(0,\frac{1}{2})\)。

Massive Dirac 粒子满足 \((i\slashed\partial-m)\psi=0\)，有 \(2s+1=2\) 个自旋态。量子场同时包含反粒子模式，因此 mode expansion 有 2 个粒子自旋态加 2 个反粒子自旋态。

Massless Weyl 费米子：左手 \(\chi_\alpha\) 和右手 \(\bar\eta^{\dot\alpha}\) 独立。在 massless 极限中 chirality 与 helicity 紧密对应：
\begin{equation}
\text{left-handed}\Leftrightarrow h=-\frac{1}{2},\qquad \text{right-handed}\Leftrightarrow h=+\frac{1}{2}.
\end{equation}
这是因为 massless 极限下 \([\gamma^5,\slashed p]=0\)（利用 \(p^2=0\)），使左右 chirality sector 解耦，各自独立描述不同螺旋度的粒子。

质量项 \(m\bar\psi\psi=m(\chi_\alpha\chi^\alpha+\bar\eta_{\dot\alpha}\bar\eta^{\dot\alpha})\) 把 \((\frac{1}{2},0)\) 与 \((0,\frac{1}{2})\) 耦合起来。因此 massive Dirac 粒子不再是单一 helicity 本征态——helicity 不再是完整 Lorentz invariant spin label。这是中微子物理中"Dirac 质量 vs Majorana 质量"讨论的表示论根源。

量子场论应用：quark、charged lepton、neutrino（Dirac 或 Majorana）、chiral gauge theory 中的 Weyl 费米子、supersymmetry 中的 gaugino/higgsino、Majorana dark matter。矢量流 \(j^\mu=\bar\psi\gamma^\mu\psi\) 与 spin-1 gauge field 的基本耦合 \(\mathcal{L}_{\text{int}}=-gA_\mu j^\mu\) 可写成协变导数 \(D_\mu=\partial_\mu+igA_\mu\) 形式 \(\mathcal{L}=\bar\psi(i\gamma^\mu D_\mu-m)\psi\)。

\section{自旋 1：photon、gluon 与 Proca 场}

四矢量场 \(A_\mu\in(\frac{1}{2},\frac{1}{2})\) 有 4 个 Lorentz 分量，但这不等于 4 个物理极化。场方程和规范对称性会把 Lorentz 分量投影到物理态。

接下来考察Massive spin-1：Proca 场。

Proca 作用量
\begin{equation}
\mathcal{L}=-\frac{1}{4}F_{\mu\nu}F^{\mu\nu}-\frac{1}{2}m^2 A_\mu A^\mu
\end{equation}
给出运动方程
\begin{equation}
\partial_\mu F^{\mu\nu}-m^2 A^\nu=0.
\end{equation}

\begin{derivation}[Proca 场的横向约束与自由度计数]
\emph{思路：}对 Proca 运动方程取散度，利用 \(F^{\mu\nu}\) 的反对称性和偏导数交换性证明 \(\partial_\nu\partial_\mu F^{\mu\nu}\equiv0\)，从而利用 \(m\neq0\) 推出 Lorenz 条件 \(\partial_\mu A^\mu=0\)。这把 4 个 Lorentz 分量约束为 3 个独立分量，与 massive spin-1 的 Wigner 分类一致。

\emph{第一步：取散度。}Proca 运动方程为 \(\partial_\mu F^{\mu\nu}-m^2 A^\nu=0\)。对两边取 \(\partial_\nu\)（协变散度）：
\begin{equation*}
\partial_\nu\partial_\mu F^{\mu\nu}-m^2\partial_\nu A^\nu=0.
\end{equation*}
目标是证明第一项恒为零，从而利用 \(m\neq0\) 得到约束。

\emph{第二步：证明 \(\partial_\nu\partial_\mu F^{\mu\nu}\equiv0\)。}展开 \(F^{\mu\nu}=\partial^\mu A^\nu-\partial^\nu A^\mu\)：
\begin{equation*}
\partial_\nu\partial_\mu F^{\mu\nu}=\partial_\nu\partial_\mu(\partial^\mu A^\nu-\partial^\nu A^\mu)=\partial_\nu\Box A^\nu-\partial_\nu\partial^\nu\partial_\mu A^\mu.
\end{equation*}
第一项 \(\partial_\nu\Box A^\nu=\Box\partial_\nu A^\nu\)（\(\Box=\partial^\mu\partial_\mu\) 与 \(\partial_\nu\) 交换，因为偏导数交换）。第二项 \(\partial_\nu\partial^\nu\partial_\mu A^\mu=\Box\partial_\mu A^\mu\)。两项相同，相减为零：
\begin{equation*}
\partial_\nu\partial_\mu F^{\mu\nu}=\Box\partial_\nu A^\nu-\Box\partial_\mu A^\nu=0.
\end{equation*}
更深刻的原因：\(\partial_\nu\partial_\mu\) 对指标 \(\mu,\nu\) 对称（偏导数交换），而 \(F^{\mu\nu}\) 对指标 \(\mu,\nu\) 反对称。对称张量与反对称张量缩并恒为零：
\begin{equation*}
\partial_\nu\partial_\mu F^{\mu\nu}=\frac{1}{2}(\partial_\nu\partial_\mu+\partial_\mu\partial_\nu)F^{\mu\nu}=\frac{1}{2}(\partial_\nu\partial_\mu F^{\mu\nu}+\partial_\mu\partial_\nu F^{\mu\nu})=0,
\end{equation*}
第二项交换哑指标 \(\mu\leftrightarrow\nu\) 得 \(\partial_\nu\partial_\mu F^{\nu\mu}=-\partial_\nu\partial_\mu F^{\mu\nu}\)（利用 \(F^{\nu\mu}=-F^{\mu\nu}\)），与第一项相加为零。

\emph{第三步：读出 Lorenz 约束。}第一项为零后，散度方程变为 \(-m^2\partial_\nu A^\nu=0\)。因 \(m\neq0\)（massive Proca 场），得 Lorenz 条件
\begin{equation*}
\partial_\mu A^\mu=0.
\end{equation*}
这个约束是运动方程的推论，而非额外的规范选择——massive 矢量场自动满足横向条件。

\emph{第四步：自由度计数。}Proca 场 \(A^\mu\) 有 4 个 Lorentz 分量。Lorenz 条件 \(\partial_\mu A^\mu=0\) 是 1 个标量约束（对每个动量模式），把 4 个分量约束为 3 个独立分量。这 \(3=2s+1\)（\(s=1\)）个极化对应 \(\lambda=-1,0,+1\)——与 massive spin-1 的 Wigner 分类一致：little group \(SO(3)\) 的 spin-1 表示有 \(2\cdot1+1=3\) 个态。物理上，\(\lambda=\pm1\) 是横向极化（自旋指向运动方向），\(\lambda=0\) 是纵向极化（自旋垂直于运动方向）——后者是 massive 矢量场特有的，无质量极限下被规范对称性消除。

\emph{第五步：与 massless 情形对比。}当 \(m=0\) 时，上述推导失效——第三步中 \(m^2=0\) 无法推出 Lorenz 条件。此时运动方程 \(\partial_\mu F^{\mu\nu}=0\) 具有 gauge redundancy \(A_\mu\to A_\mu+\partial_\mu\alpha\)，比 Lorenz 条件更强：不仅消除纵向分量，还消除残余规范自由度。最终只留 \(h=\pm1\) 两个横向 helicity。massive \(\to\) massless 的极限使自由度从 3 跳到 2，纵向极化 \(\lambda=0\) 被 gauge 吸收——这就是电弱对称破缺中 \(W/Z\) 玻色子获得质量时"吃掉"Goldstone boson 的运动学机制。

\emph{Stueckelberg 机制。}Proca 理论可嵌入一个形式上 gauge-invariant 的框架：引入标量场 \(\phi\)（Stueckelberg 场），将 Proca Lagrangian 改写为
\begin{equation*}
\mathcal{L}_{\mathrm{St}}=-\frac{1}{4}F_{\mu\nu}F^{\mu\nu}+\frac{1}{2}m^2(A_\mu-\partial_\mu\phi)(A^\mu-\partial^\mu\phi),
\end{equation*}
此 Lagrangian 在 \(A_\mu\to A_\mu+\partial_\mu\alpha,\;\phi\to\phi+\alpha\) 下不变。取 unitary gauge \(\phi=0\) 恢复 Proca 形式；取 \(R_\xi\) gauge 可得具有良好 UV 行为的传播子。Stueckelberg 场 \(\phi\) 恰好提供被 gauge 吸收的纵向自由度——在自发破缺中，\(\phi\) 就是 Goldstone boson，"被吃掉"后矢量场获得质量。

\emph{Proca 传播子。}在动量空间，Proca 运动方程 \((g_{\mu\nu}\Box-\partial_\mu\partial_\nu+m^2g_{\mu\nu})A^\nu=0\) 的逆给出 Feynman 传播子
\begin{equation*}
D_{\mu\nu}(p)=\frac{-i}{p^2-m^2+i\epsilon}\left(g_{\mu\nu}-\frac{p_\mu p_\nu}{m^2}\right).
\end{equation*}
括号内正是极化完备关系 \(\sum_\lambda\epsilon_\mu^{(\lambda)}\epsilon_\nu^{(\lambda)*}\)，确认传播子只投影到 3 个物理极化子空间。与 massless 光子传播子 \(\frac{-ig_{\mu\nu}}{p^2+i\epsilon}\)（Lorenz gauge）对比，Proca 传播子的 \(p_\mu p_\nu/m^2\) 项在高能 \(p^2\gg m^2\) 时趋于常数，不随 \(p^2\) 增长——这是 massive 矢量场可重整化的关键。但单个 Proca 场与费米子 Yukawa 耦合会破坏可重整性；只有通过自发破缺（Higgs 机制）嵌入非阿贝尔规范理论，才能同时保持规范对称性和可重整性。

\emph{电弱破缺中的实现。}标准模型中 \(W^\pm\) 和 \(Z\) 玻色子是 Proca 场的物理实例。Higgs 双重态 \(\Phi\) 获得真空期望值后，规范破缺 \(SU(2)_L\times U(1)_Y\to U(1)_{\mathrm{em}}\) 使 3 个规范玻色子获得质量 \(m_W=\frac{gv}{2}\)、\(m_Z=\frac{gv}{2\cos\theta_W}\)，对应的 3 个 Goldstone 模式被 Stueckelberg 机制吸收为纵向极化。剩余的 photon 保持无质量。\(W/Z\) 的纵向极化态在高能极限下的散射振幅涉及等效定理：\(A(W_L W_L\to W_L W_L)\simeq A(\phi\phi\to\phi\phi)+\mathcal{O}(m_W/E)\)，其中 \(\phi\) 是被吃掉的 Goldstone——这保证纵向 \(W\) 散射在 TeV 尺度不破坏幺正性，是 Higgs 粒子存在的理论信号。

\emph{极化矢量的显式构造。}在静止系 \(p^\mu=(m,\mathbf{0})\)，三个极化矢量取纯空间形式 \(\epsilon_\mu^{(\lambda)}=(0,\mathbf{e}^{(\lambda)})\)，其中 \(\mathbf{e}^{(\lambda)}\) 是 \(SO(3)\) 的 spin-1 基矢。对一般动量 \(p^\mu=(E,0,0,p)\)（沿 \(z\) 轴），Lorentz boost 给出
\begin{equation*}
\epsilon^{(\pm1)}=\frac{1}{\sqrt{2}}(0,1,\pm i,0),\qquad\epsilon^{(0)}=\frac{1}{m}(p,0,0,E),
\end{equation*}
可直接验证 \(p^\mu\epsilon_\mu^{(\lambda)}=0\)（横向条件）及正交完备 \(\epsilon^{(\lambda)*}\cdot\epsilon^{(\lambda')}=-\delta_{\lambda\lambda'}\)。纵向极化 \(\epsilon^{(0)}\) 在高能 \(E\gg m\) 时 \(\epsilon_\mu^{(0)}\sim p_\mu/m+\mathcal{O}(m/E)\)，主导项 \(p_\mu/m\) 正是等效定理中 Goldstone-矢量对应的运动学根源。

\emph{\(R_\xi\) 规范族。}Stueckelberg 理论的 \(R_\xi\) gauge-fixing 项 \(\mathcal{L}_{\mathrm{gf}}=-\frac{1}{2\xi}(\partial_\mu A^\mu-\xi m\phi)^2\) 给出一族传播子
\begin{equation*}
D_{\mu\nu}^{(\xi)}(p)=\frac{-i}{p^2-m^2+i\epsilon}\left(g_{\mu\nu}-(1-\xi)\frac{p_\mu p_\nu}{p^2-\xi m^2}\right),
\end{equation*}
\(\xi=1\)（Feynman gauge）给出 \(D_{\mu\nu}=-ig_{\mu\nu}/(p^2-m^2)\)——形式上与标量传播子相同，计算最简；\(\xi\to\infty\)（unitary gauge）恢复 Proca 传播子。不同 \(\xi\) 下物理振幅相同（Ward--Takahashi 恒等式保证），但中间步骤的 Feynman 图复杂度差异显著——这体现了规范冗余在计算中的实用价值。
\end{derivation}

物理极化满足 \(p^\mu\epsilon_\mu^{(\lambda)}=0\)，completeness relation 为
\begin{equation}\label{eq:ms-curved-pol-sum}
\boxed{
\sum_{\lambda=-1}^{1}\epsilon_\mu^{(\lambda)}(p)\,\epsilon_\nu^{(\lambda)\ast}(p)=g_{\mu\nu}+\frac{p_\mu p_\nu}{m^2}.
}
\end{equation}
右端第一项是 4 维单位矩阵（4 个自由度），第二项 \(\frac{p_\mu p_\nu}{m^2}\) 减去纵向极化（1 个非物理自由度），留下 3 个物理极化。

接下来考察Massless spin-1：Maxwell 场。

Maxwell 理论 \(\mathcal{L}=-\frac{1}{4}F_{\mu\nu}F^{\mu\nu}\) 具有 gauge redundancy \(A_\mu\to A_\mu+\partial_\mu\alpha\)。这个规范对称性比 Proca 的约束更强：不仅消除纵向极化（\(\partial_\mu A^\mu=0\)，Lorenz gauge），还消除剩余的 gauge freedom（残余规范变换 \(\alpha\) 满足 \(\Box\alpha=0\)）。最终只剩 \(h=\pm1\) 两个 transverse helicity——与 massless spin-1 的 Wigner 分类一致：little group \(ISO(2)\) 的 finite helicity 表示，parity 配对后维数 2。

量子场论应用：photon（\(U(1)\) gauge boson，\(h=\pm1\)）、gluon（\(SU(3)_c\) adjoint，仍为 \(h=\pm1\)，color index 是独立的内部自由度）、\(W^\pm/Z\)（电弱破缺后 massive，\(\lambda=-1,0,+1\)，纵向极化与 Higgs/Goldstone 相关——Goldstone equivalence theorem 将高能极限下 longitudinal \(W/Z\) 与相应 Goldstone mode 联系起来）。

\section{自旋 3/2：Rarita--Schwinger 与 gravitino}

vector-spinor \(\psi_\mu\) 同时具有矢量与旋量指标。在 Lorentz 表示中，\((\frac{1}{2},\frac{1}{2})\otimes\big((\frac{1}{2},0)\oplus(0,\frac{1}{2})\big)\) 不仅包含 spin-\(\frac{3}{2}\) 还包含 spin-\(\frac{1}{2}\) 成分。必须通过方程与约束去除低自旋部分。

Massive spin-\(\frac{3}{2}\) 需要三个条件共同约束：
\begin{equation}
(\slashed p-m)\psi_\mu=0,\qquad p^\mu\psi_\mu=0,\qquad \gamma^\mu\psi_\mu=0.
\end{equation}
第一个是 Dirac 方程（确定质量），第二个是横向约束（去除矢量方向的冗余），第三个是 gamma-traceless 约束（去除 spin-\(\frac{1}{2}\) 成分）。共同留下 \(2s+1=4\) 个物理态 \(\lambda=-\frac{3}{2},-\frac{1}{2},+\frac{1}{2},+\frac{3}{2}\)。

Massless Rarita--Schwinger 场具有 gauge symmetry \(\delta\psi_\mu=\partial_\mu\epsilon\)（\(\epsilon\) 是旋量规范参数）。gauge redundancy 消除低 helicity 成分后保留 \(h=\pm\frac{3}{2}\)。在曲率时空中，gauge symmetry 会受到 curvature obstruction：一般背景上 \([\nabla_\mu,\nabla_\nu]\neq0\) 破坏 gauge 一致性，通常需要 \(R_{\mu\nu}=0\)（Ricci-flat）等背景限制。

量子场论应用：超引力中的 gravitino（super-Higgs mechanism 后可获得质量）、强子 \(\Delta\)-type resonance 的有效描述。

\section{自旋 2：graviton 与 massive spin-2}

对称二阶张量 \(h_{\mu\nu}=h_{\nu\mu}\) 在四维有 10 个分量。与 spin-1 类似，需要区分 massive 和 massless。

接下来考察Massless spin-2：graviton。

线性化 diffeomorphism \(\delta h_{\mu\nu}=\partial_\mu\xi_\nu+\partial_\nu\xi_\mu\)（\(\xi^\mu\) 是 4 个规范参数）的规范冗余与约束最终只留下 \(h=\pm2\) 两个物理 helicity。自由度计数：10 个分量 \(-\) 4 个规范 \(-\) 4 个约束（来自场方程的 Bianchi 结构）\(=2\)。

spin-2 场与能动量张量的普适耦合
\begin{equation}\label{eq:ms-curved-spin2-coupling}
\boxed{\mathcal{L}_{\text{int}}=-\frac{\kappa}{2}h_{\mu\nu}T^{\mu\nu}.}
\end{equation}
这是量子场论中重要的结构类比：spin-1 gauge boson couples to conserved current \(J^\mu\)（\(\partial_\mu J^\mu=0\)），spin-2 graviton couples to conserved stress tensor \(T^{\mu\nu}\)（\(\partial_\mu T^{\mu\nu}=0\)）。耦合源的守恒性正是 gauge invariance 的推论——规范对称性要求源无散度。

接下来考察Massive spin-2：Fierz--Pauli。

Fierz--Pauli 理论中 massive spin-2 有 \(2s+1=5\) 个极化 \(\lambda=-2,-1,0,+1,+2\)。可用 massive spin-1 极化矢量构造极化张量：
\begin{equation}
\epsilon_{\mu\nu}^{(\pm 2)}=\epsilon_\mu^{(\pm)}\epsilon_\nu^{(\pm)},\qquad
\epsilon_{\mu\nu}^{(\pm 1)}=\frac{1}{\sqrt{2}}(\epsilon_\mu^{(\pm)}\epsilon_\nu^{(0)}+\epsilon_\mu^{(0)}\epsilon_\nu^{(\pm)}),
\end{equation}
helicity-0 取适当的 traceless combination。massive \(\to\) massless 极限在相互作用理论中并不平凡——helicity-0 mode 导致著名的 Vainshtein mechanism 和 strong coupling scale 问题。

\section{一般自旋的 Fronsdal 场}

接下来考察Massive integer spin-\(s\)。

用 totally symmetric tensor \(\phi_{\mu_1\cdots\mu_s}=\phi_{(\mu_1\cdots\mu_s)}\) 描述。Fierz--Pauli 型条件
\begin{equation}
(\Box-m^2)\phi_{\mu_1\cdots\mu_s}=0,\qquad \partial^{\mu_1}\phi_{\mu_1\cdots}=0,\qquad \phi^\rho{}_{\rho\cdots}=0
\end{equation}
（Klein--Gordon + transverse + traceless）把 Lorentz tensor 中多余的低 spin 成分投影掉，留下 \(2s+1\) 个 massive spin states。

接下来考察Massless integer spin-\(s\)：Fronsdal 场。

Fronsdal gauge transformation
\begin{equation}
\delta\phi_{\mu_1\cdots\mu_s}=s\,\partial_{(\mu_1}\epsilon_{\mu_2\cdots\mu_s)},\qquad \epsilon^\rho{}_{\rho\cdots}=0
\end{equation}
（规范参数 traceless）。对 \(s\ge4\)，Fronsdal field 要求 double-traceless \(\phi^{\rho\sigma}{}_{\rho\sigma\cdots}=0\)。Fronsdal 方程 \(\mathcal{F}_{\mu_1\cdots\mu_s}=0\)，其中 Fronsdal 算符
\begin{equation}
\mathcal{F}_{\mu_1\cdots\mu_s}=\Box\phi_{\mu_1\cdots\mu_s}-s\,\partial_{(\mu_1}(\partial\cdot\phi)_{\mu_2\cdots\mu_s)}+\frac{s(s-1)}{2}\partial_{(\mu_1}\partial_{\mu_2}\phi'_{\mu_3\cdots\mu_s)},
\end{equation}
在四维传播 \(h=\pm s\) 两个物理 helicity。

接下来考察为什么 gauge symmetry 与 helicity 直接相关。

高 rank Lorentz field 往往包含多个 spin/helicity sector。对 massless field，gauge symmetry 的核心作用是消除非物理 polarizations：
\begin{equation}
A_\mu\Rightarrow h=\pm 1,\qquad \psi_\mu\Rightarrow h=\pm\frac{3}{2},\qquad h_{\mu\nu}\Rightarrow h=\pm 2.
\end{equation}
物理上，gauge symmetry 正是 little group \(ISO(2)\) 平移子群的实现——规范变换把多余的"纵向"分量 gauge away，只留下 little group 的旋转部分（helicity）作为物理标签。因此
\begin{equation}
\text{massless gauge symmetry}\Longrightarrow\text{physical helicity projection}.
\end{equation}

\section{螺旋度振幅与 little-group scaling}

对 massless particle，四维动量可以写为 spinor 双线性
\begin{equation}
p_{\alpha\dot\alpha}=\lambda_\alpha\tilde\lambda_{\dot\alpha},
\end{equation}
其中 \(\lambda_\alpha\) 和 \(\tilde\lambda_{\dot\alpha}\) 是两分量旋量。little-group scaling \(\lambda\to t\lambda\)、\(\tilde\lambda\to t^{-1}\tilde\lambda\)（保持 \(p_{\alpha\dot\alpha}\) 不变）要求 helicity-\(h\) 外线态的振幅满足齐次性条件
\begin{equation}\label{eq:ms-curved-little-group}
\boxed{\mathcal{A}_n\to t_i^{-2h_i}\mathcal{A}_n}
\end{equation}
对每条外线 \(i\) 分别成立。这把 helicity 表示论直接转化为 scattering amplitude 的齐次性约束——振幅的 little-group scaling 完全确定了它的 spinor 结构。

最基本的应用：Yang--Mills 三点振幅
\begin{equation}
A_3(1^-,2^-,3^+)=g\,\frac{\langle 12\rangle^3}{\langle 23\rangle\langle 31\rangle},
\end{equation}
gravity 三点振幅
\begin{equation}
M_3(1^{-2},2^{-2},3^{+2})=\kappa\,\frac{\langle 12\rangle^6}{\langle 23\rangle^2\langle 31\rangle^2}.
\end{equation}
gravity 振幅在结构上是 Yang--Mills 振幅的"平方"——这是 Kawai--Lewellen--Tye (KLT) double-copy 结构的最简单信号，也是现代量子引力中"gravity = (gauge theory)\(^2\)"猜想的振幅级证据。

\section{从 Lorentz 表示到物理态的统一框架}

把所有场类型并列，可以看清"表示 \(\to\) 物理态"的统一逻辑：

\begin{center}
\begin{tabular}{llll}
\hline
场 & Lorentz 表示 & 约束/gauge & 物理态 \\
\hline
\(\phi\) & scalar & 无 & \(s=0\) \\
\(\psi\) massive & Dirac & Dirac eq. & \(s=\frac{1}{2}\), 2 states \\
\(\chi\) massless & Weyl & Weyl eq. & 单 chirality helicity \\
\(A_\mu\) massive & vector & \(p\cdot A=0\) & \(s=1\), 3 states \\
\(A_\mu\) massless & gauge vector & \(A\sim A+\partial\alpha\) & \(h=\pm 1\) \\
\(\psi_\mu\) massive & vector-spinor & \(p\cdot\psi=0,\gamma\cdot\psi=0\) & \(s=\frac{3}{2}\), 4 states \\
\(\psi_\mu\) massless & gauge vector-spinor & \(\psi\sim\psi+\partial\epsilon\) & \(h=\pm\frac{3}{2}\) \\
\(h_{\mu\nu}\) massive & symmetric tensor & transverse + traceless & \(s=2\), 5 states \\
\(h_{\mu\nu}\) massless & gauge symmetric tensor & linear diffeo. & \(h=\pm 2\) \\
\hline
\end{tabular}
\end{center}

真正的逻辑不是"张量的 rank = 自旋"——这只是对 massive free field 的巧合。一般地，
\begin{equation}\label{eq:ms-curved-unified-logic}
\boxed{
\text{Lorentz 表示}+\text{场方程}+\text{约束/gauge}\Longrightarrow\text{物理 Poincaré spin/helicity}.
}
\end{equation}
三层结构：第一层是 Lorentz 表示 \(\Sigma^{\mu\nu}\)（决定场如何变换）；第二层是 Poincaré Casimir \(P^2,W^2\)（决定质量和自旋/helicity）；第三层是场方程+约束+gauge（把大的 Lorentz 表示投影到物理 Hilbert space）。理解标量、光子、Dirac 费米子、graviton 和高自旋场的统一框架，就是理解这三层结构如何协同工作。

\section*{本章小结}

曲率时空中，统一的协变导数 \(D_\mu=\partial_\mu-\frac{i}{2}\omega_{\mu ab}\Sigma^{ab}\) 把平直时空的所有场论结构一并推广，不同表示只是 \(\Sigma^{ab}\) 不同——标量为零、矢量为 \(d\times d\) 矩阵、旋量为 Clifford 代数构造。diffeomorphism invariance 自动给出协变守恒 \(\nabla_\mu T^{\mu\nu}=0\)，Killing 矢量把它转化为具体的守恒荷。当度规本身成为动力学变量，Einstein--Hilbert 作用量给出 Einstein 方程，Bianchi 恒等式保证物质守恒——几何方程的恒等式与物质守恒律自洽。

Wigner 分类把场的 Lorentz 表示与物理单粒子的 Poincaré 不可约表示联系起来：massive 粒子由 \(SO(3)\) little group 分类，有 \(2s+1\) 个极化；massless 粒子由 \(ISO(2)\) little group 分类，gauge symmetry 把多余的 Lorentz 分量投影到 \(h=\pm s\) 两个物理 helicity。整个链条可以压缩为：
\begin{equation*}
\text{Lie group}\to\text{Lie algebra}\to\text{representation}\to\text{field transformation}\to\text{action}\to\text{Euler--Lagrange}\to\text{Noether current}\to\text{conserved charge}.
\end{equation*}
平直时空中 \(P^\mu,M^{\mu\nu}\) 是 Poincaré 生成元；曲率时空中全局 Poincaré 对称被 diffeomorphism covariance + 局域 Lorentz 对称取代，但协变导数的统一形式使所有场类型的推导保持一致。螺旋度振幅的 little-group scaling 把表示论直接转化为散射振幅的齐次性约束，是现代量子场论中表示论最直接的应用。
